%% ============================================================
%%  Article scientifique — double colonne IEEE
%%  Intégration de nouvelles informations dans les LLMs :
%%  comparaison de sept stratégies sur corpus français hétérogène
%%  Compilez avec : pdflatex article_scientifique.tex (x2)
%% ============================================================

\documentclass[10pt,journal,compsoc]{IEEEtran}

%% ---- Packages ----
\usepackage[utf8]{inputenc}
\usepackage[T1]{fontenc}
\usepackage[french]{babel}
\usepackage{amsmath,amssymb}
\usepackage{booktabs}
\usepackage{multirow}
\usepackage{graphicx}
\usepackage{xcolor}
\usepackage{hyperref}
\usepackage{microtype}
\usepackage{array}
\usepackage{tabularx}
\usepackage{caption}
\usepackage{subcaption}
\usepackage{url}
\usepackage{cite}
\usepackage{float}

%% ---- Couleurs personnalisées ----
\definecolor{best}{RGB}{0,128,0}       % vert : meilleur résultat
\definecolor{worst}{RGB}{180,0,0}      % rouge : moins bon résultat

%% ---- Commandes utilitaires ----
\newcommand{\best}[1]{\textbf{\textcolor{best}{#1}}}   % meilleur score
\newcommand{\worst}[1]{\textcolor{worst}{#1}}           % pire score
\newcommand{\bs}{BERTScore}
\newcommand{\rl}{ROUGE-L}
\newcommand{\met}{METEOR}
\newcommand{\fc}{\textit{function calling}}

%% ---- Métadonnées ----
\title{Intégration de nouvelles informations dans les LLMs :\\
Baseline, RAG, Fine-tuning LoRA, FT+RAG, RAFT,\\
Reranking et Function Calling sur corpus français hétérogène}

\author{%
  \IEEEauthorblockN{Abdoul Karim Coulibaly}\\
  \IEEEauthorblockA{%
    CESI École d'Ingénieurs, Saint-Nazaire, France\\
    \textit{Projet pédagogique --- FISE A4 Data Science \& IA, 2025--2026}\\
    \href{mailto:abdoulkarim.coulibaly@viaacesi.fr}{abdoulkarim.coulibaly@viaacesi.fr}\\
    \url{https://github.com/coulby45/llm-integration-pipeline}
  }
}

\markboth{Projet Article Scientifique --- CESI Saint-Nazaire, 2026}{}

\begin{document}

\maketitle

%% ============================================================
\begin{abstract}
Les grands modèles de langage (LLMs) souffrent d'une date de coupure
figée~: \textsc{LLaMA~3.1~8B}, dont la coupure s'arrête à décembre
2023, ignore toute publication ou événement postérieur. Ce verrou
constitue un frein majeur pour les applications nécessitant des
connaissances fraîches. Nous proposons une évaluation comparative
de \textbf{sept stratégies d'intégration d'information} sur ce socle
commun~: (1)~Baseline (Groq sans contexte), (2)~RAG dense (FAISS +
bi-encodeur multilingue), (3)~Fine-tuning LoRA, (4)~FT+RAG (adaptateur
LoRA combiné à une récupération FAISS à l'inférence), (5)~RAFT
(\textit{Retrieval-Augmented Fine-Tuning}), (6)~RAG + reranking
cross-encoder, et (7)~Function calling avec outil \texttt{search\_docs}.
L'évaluation est conduite sur un corpus \textbf{100\,\% français}
multi-source (Wikipedia, HAL.science, actualités, code de la route)
de 1\,450 paires question--réponse (\textit{n}$_\text{test}$\,=\,290),
selon des métriques lexicales (F1, \rl, \met), sémantiques (\bs{}
avec intervalles de confiance bootstrap à 95\,\%), de fidélité
(proxy d'hallucination) et subjectives (LLM-as-judge via
Claude Sonnet~4.6). Nos résultats montrent que le \textbf{reranking}
et le \textbf{RAG} dominent toutes les métriques qualitatives avec
respectivement un F1 de 36,5\,\% et 32,8\,\%, contre 16,3\,\% pour
la baseline, tandis que les méthodes par fine-tuning (LoRA, FT+RAG,
RAFT) \textbf{n'améliorent pas} la baseline sur ce protocole. Le
\textit{function calling} se positionne en intermédiaire (F1\,=\,23,9\,\%).
Ces résultats soulèvent une discussion sur le coût computationnel,
la fraîcheur des données et les biais méthodologiques inhérents
à l'évaluation automatique.
\end{abstract}

\begin{IEEEkeywords}
LLM, RAG, LoRA, RAFT, reranking, function calling, évaluation,
corpus français, hallucination, BERTScore, FAISS, LLM-as-judge
\end{IEEEkeywords}


%% ============================================================
\section{Introduction}
\label{sec:intro}

\subsection{Contexte et motivation}

Les LLMs comme GPT-4~\cite{openai2023gpt4},
Mistral~7B~\cite{jiang2023mistral} ou
\textsc{LLaMA}~3.1~8B~\cite{meta2024llama3} ont démontré des capacités
remarquables de compréhension et de génération de texte. Cependant,
ils partagent une limitation fondamentale~: leur connaissance du monde
est figée à leur \textbf{date de coupure d'entraînement}. Pour
\textsc{LLaMA}~3.1~8B, cette coupure s'arrête à décembre 2023~; le
modèle ignore donc les dernières publications scientifiques, les
nouvelles lois promulguées, ou les événements d'actualité récents.

Cette contrainte engendre deux problèmes pratiques majeurs~:
(i)~l'\textbf{obsolescence} des réponses sur des sujets évolutifs~;
(ii)~les \textbf{hallucinations}~\cite{ji2023hallucination}, où le
modèle génère avec confiance des informations factuellement incorrectes
faute de données récentes. La question centrale de ce travail est donc~:

\begin{quote}
  \textit{Comment intégrer efficacement de nouvelles connaissances dans
  un LLM sans le réentraîner entièrement, et quelle stratégie offre le
  meilleur compromis qualité/coût sur du français ?}
\end{quote}

\subsection{Verrou scientifique}

La littérature existante souffre de trois biais récurrents. Premièrement,
les comparaisons mobilisent souvent des \textbf{modèles différents} pour
chaque méthode, rendant les conclusions difficilement généralisables.
Deuxièmement, la quasi-totalité des benchmarks sont en
\textbf{anglais}~\cite{rajpurkar2016squad}, ce qui ne rend pas compte
des spécificités du traitement du français. Troisièmement, les
évaluations ignorent fréquemment la \textbf{nature du domaine}
(textes techniques \textit{vs} actualités) et la \textbf{complexité des
questions} (factuelle simple \textit{vs} raisonnement multi-sauts).

\subsection{Contributions}

Ce travail apporte les contributions suivantes~:

\begin{enumerate}
  \item Un \textbf{corpus 100\,\% français} multi-source et multi-domaine
        (Wikipedia, HAL.science, presse, code de la route) avec quatre
        types de domaine et cinq types de questions.
  \item Un \textbf{protocole de comparaison équitable}~: les sept méthodes
        partagent le même socle \textsc{LLaMA}~3.1~8B.
  \item Une \textbf{évaluation multi-axe} incluant métriques lexicales,
        sémantiques, proxy d'hallucination, latence, impact écologique
        et évaluation subjective LLM-as-judge.
  \item Une \textbf{analyse croisée} méthodes $\times$ domaines $\times$
        types de question, avec intervalles de confiance bootstrap.
  \item Des \textbf{recommandations pratiques} pour le déploiement en
        contexte francophone contraint.
\end{enumerate}

\subsection{Plan de l'article}

La section~\ref{sec:related} présente l'état de l'art. La
section~\ref{sec:methodo} décrit notre méthodologie et les sept
pipelines. Les résultats sont présentés en section~\ref{sec:results}
et discutés en section~\ref{sec:discussion}. La
section~\ref{sec:conclusion} conclut et ouvre des perspectives.


%% ============================================================
\section{État de l'art}
\label{sec:related}

\subsection{Grands modèles de langage}

Depuis l'introduction du mécanisme d'attention par \textit{Transformer}
\cite{vaswani2017attention}, les LLMs ont connu une croissance
exponentielle de leurs capacités~: de GPT-2~\cite{radford2019language}
à GPT-3~\cite{brown2020language} puis à
\textsc{LLaMA}~2~\cite{touvron2023llama} et
\textsc{LLaMA}~3~\cite{meta2024llama3}. Notre socle expérimental,
\textsc{LLaMA}~3.1~8B, dispose d'une fenêtre de contexte de 128k tokens
et d'une quantification 4-bit via Unsloth~\cite{unsloth2024}, ce qui
le rend utilisable sur GPU grand public (NVIDIA A100 dans nos
expériences sur Google Colab).

\subsection{Problème de la mise à jour des connaissances}

Le \textit{catastrophic forgetting}~\cite{mccloskey1989catastrophic,
kirkpatrick2017overcoming} décrit la tendance d'un réseau de neurones
à oublier des connaissances antérieures lors d'un nouvel entraînement.
Ji~et~al.~\cite{ji2023hallucination} ont formalisé le problème des
hallucinations dans les LLMs. Des approches d'édition de connaissances
telles que ROME~\cite{meng2022locating} ou
SERAC~\cite{mitchell2022memory} permettent de modifier des faits
spécifiques sans réentraînement complet, mais peinent à généraliser.

\subsection{Génération augmentée par récupération (RAG)}

Lewis~et~al.~\cite{lewis2020rag} ont introduit le paradigme RAG~:
au moment de l'inférence, un retrieveur dense récupère des passages
pertinents depuis une base documentaire, qui sont ensuite injectés
dans le prompt du générateur. Cette approche découple la mémorisation
(encodée dans l'index) de la génération (produite par le LLM),
permettant de mettre à jour les connaissances par simple réindexation.

L'indexation vectorielle efficace repose sur FAISS~\cite{johnson2019faiss}.
Les représentations sémantiques sont obtenues via des
Sentence-Transformers~\cite{reimers2019sentencebert}, dont les variantes
multilingues couvrent le français. Le reranking
cross-encoder~\cite{nogueira2019passage} affine le classement initial
en évaluant chaque paire (requête, passage) conjointement.

\subsection{Fine-tuning efficace --- LoRA et RAFT}

Le fine-tuning complet de LLMs est prohibitif en ressources.
LoRA~\cite{hu2022lora} résout ce problème en décomposant les mises à
jour de poids en matrices de bas rang~: si $\Delta W = AB$ avec
$A \in \mathbb{R}^{d \times r}$ et $B \in \mathbb{R}^{r \times d}$,
le nombre de paramètres passe de $d^2$ à $2rd$ avec $r \ll d$.
QLoRA~\cite{dettmers2023qlora} ajoute une quantification 4-bit,
réduisant encore l'empreinte mémoire.

RAFT (\textit{Retrieval-Augmented Fine-Tuning})~\cite{hong2024raft}
propose d'entraîner le modèle avec des contextes mixtes~: des passages
pertinents (\textit{oracle documents}) mélangés à des distracteurs
(\textit{distractor documents}). Le modèle apprend ainsi à ignorer
le bruit retrieval dès l'entraînement, réduisant le désalignement
entre phase d'entraînement (avec récupération) et inférence.

\subsection{Function calling et agents à outils}

Le \textit{function calling} permet au LLM de décider dynamiquement
quand et comment interroger un outil externe~\cite{schick2023toolformer}.
Plutôt que d'injecter systématiquement des passages dans le prompt, le
modèle émet un appel d'outil structuré (JSON), l'exécution restant côté
client. Cette approche est particulièrement adaptée aux flux agentiques
nécessitant des politiques de recherche conditionnelle.

\subsection{Évaluation des LLMs}

L'évaluation repose classiquement sur des métriques lexicales~:
F1 token-level~\cite{rajpurkar2016squad}, ROUGE-L~\cite{lin2004rouge},
et METEOR~\cite{banerjee2005meteor}. Ces métriques pénalisent les
paraphrases correctes et ne captent pas la cohérence factuelle.
BERTScore~\cite{zhang2020bertscore} remédie partiellement à ce problème
en calculant la similarité sémantique via des représentations
contextualisées. Des frameworks comme RAGAS~\cite{es2024ragas} ont
été développés pour évaluer spécifiquement les systèmes RAG. Plus
récemment, les approches \textit{LLM-as-judge} utilisent un modèle
de grande taille comme évaluateur automatique, offrant une mesure
qualitative complémentaire aux métriques lexicales.


%% ============================================================
\section{Méthodologie}
\label{sec:methodo}

\subsection{Architecture générale du pipeline}

L'ensemble des sept méthodes comparées s'appuie sur le même socle~:
\textbf{\textsc{LLaMA}~3.1~8B}. Les méthodes sans adaptation des poids
(Baseline, RAG, Reranking, Function calling) utilisent le modèle via
l'API Groq (\texttt{llama-3.1-8b-instant}). Les méthodes avec
adaptation (Fine-tuning LoRA, FT+RAG, RAFT) mobilisent la version
quantifiée 4-bit (\texttt{unsloth/Meta-Llama-3.1-8B-bnb-4bit}) sur
GPU NVIDIA A100 via Google Colab.

%%
%% [FIGURE 1 — À INSÉRER]
%% Architecture globale du pipeline : corpus → génération Q&A (Groq)
%% → 7 méthodes → évaluation (09_evaluation.ipynb)
%% Fichier : figures/fig1_architecture_globale.pdf
%%
\begin{figure}[H]
  \centering
  \fbox{\parbox{0.9\linewidth}{\centering
    \textit{[Figure 1 --- Architecture globale du pipeline]}\\[4pt]
    Corpus (4 sources) $\rightarrow$ Q\&A Claude Sonnet~4.6 $\rightarrow$
    7 méthodes $\rightarrow$ Évaluation \texttt{09\_evaluation.ipynb}
  }}
  \caption{Vue d'ensemble du pipeline expérimental. Les sept méthodes
           partagent le même socle \textsc{LLaMA}~3.1~8B.}
  \label{fig:pipeline}
\end{figure}

\subsection{Constitution du corpus}
\label{subsec:corpus}

Nous constituons un corpus \textbf{100\,\% français} multi-source
(Table~\ref{tab:corpus}). Les questions et réponses de référence sont
générées par Claude Sonnet~4.6 via l'API Anthropic, avec un contexte
(\texttt{context}) extrait verbatim du document source. Le générateur
du dataset étant \textbf{distinct} du modèle d'inférence testé
(\textsc{LLaMA}~3.1~8B), aucun biais circulaire n'affecte l'évaluation.

\begin{table}[h]
  \centering
  \caption{Statistiques du corpus expérimental (split 80/20, seed\,=\,42,
           stratification par \texttt{recency\_category})}
  \label{tab:corpus}
  \renewcommand{\arraystretch}{1.25}
  \begin{tabular}{lcccc}
    \toprule
    \textbf{Dataset} & \textbf{Source} & \textbf{Train} &
    \textbf{Test} & \textbf{Q/doc} \\
    \midrule
    Technique   & Wikipedia FR      & 320 & 80  & 5 \\
    Multi-sauts & HAL.science FR    & 360 & 90  & 5 \\
    Temporel    & Le Monde / FI     & 400 & 100 & 5 \\
    Juridique   & Code de la route  & 80  & 20  & 5 \\
    \midrule
    \multicolumn{2}{l}{\textbf{Total}} &
    \textbf{1\,160} & \textbf{290} & — \\
    \bottomrule
  \end{tabular}
\end{table}

La catégorie de récence est déduite de la date ISO du document~:
\texttt{récent} ($\geq$\,2024), \texttt{intermédiaire} (2022--2023),
\texttt{fondamental} ($<$\,2022). Les types de questions varient~:
factuelle, synthèse, compréhension (Wikipedia, presse) et
simple/complexe multi-sauts (HAL). Le domaine juridique est constitué
de segments du PDF officiel du code de la route ($\approx$\,1\,400 mots
par segment), indexés dans le même FAISS que les autres sources.

\subsection{Méthode 1 --- Baseline}

La baseline consiste à interroger \textsc{LLaMA}~3.1~8B via Groq
\textbf{sans aucun contexte externe}. Le prompt minimal en français
guide le modèle à répondre depuis ses paramètres figés. Cette méthode
représente la limite inférieure~: les connaissances antérieures à
décembre 2023 accessibles au modèle sans aide documentaire.

\subsection{Méthode 2 --- RAG dense}

%%
%% [FIGURE 2 — À INSÉRER]
%% Pipeline RAG dense (flowchart TB depuis schemas_methodes.md)
%% Fichier : figures/fig2_rag_pipeline.pdf
%%
\begin{figure}[H]
  \centering
  \fbox{\parbox{0.9\linewidth}{\centering
    \textit{[Figure 2 --- Pipeline RAG dense]}\\[4pt]
    Question $\rightarrow$ Bi-encodeur MiniLM $\rightarrow$
    FAISS (top-5) $\rightarrow$ Prompt + passages $\rightarrow$ Groq
  }}
  \caption{Pipeline RAG dense. Les documents sont indexés en chunks
           de 200 mots (chevauchement 50 mots) dans un index FAISS
           \texttt{IndexFlatIP}.}
  \label{fig:rag}
\end{figure}

L'embedding est assuré par
\texttt{paraphrase-multilingual-MiniLM-L12-v2} (dimension 384,
FAISS \texttt{IndexFlatIP} sur vecteurs normalisés L2). Le corpus est
découpé en chunks de 200 mots avec chevauchement de 50 mots~; ce
paramétrage a été préféré à des fenêtres courtes (50--100 mots) qui
dégradaient les performances du RAG sous la baseline dans une
configuration préliminaire. Les 5 passages les plus proches
(titre + texte) sont injectés dans le prompt Groq.

\subsection{Méthode 3 --- Fine-tuning LoRA}

%%
%% [FIGURE 3 — À INSÉRER]
%% Pipeline Fine-tuning LoRA (flowchart TB depuis schemas_methodes.md)
%% Fichier : figures/fig3_lora_pipeline.pdf
%%
\begin{figure}[H]
  \centering
  \fbox{\parbox{0.9\linewidth}{\centering
    \textit{[Figure 3 --- Pipeline Fine-tuning LoRA]}\\[4pt]
    train.json $\rightarrow$ Unsloth + LoRA (rank 32) $\rightarrow$
    Adaptateur $\rightarrow$ LLaMA 8B + adaptateur $\rightarrow$
    Prédiction locale
  }}
  \caption{Fine-tuning LoRA via Unsloth sur GPU A100 (Google Colab).
           L'entraînement utilise le gabarit Alpaca (format SFT).}
  \label{fig:lora}
\end{figure}

Le fine-tuning s'appuie sur Unsloth~\cite{unsloth2024} avec la
configuration LoRA suivante~: rank\,=\,32, $\alpha$\,=\,32,
dropout\,=\,0.05, 7 modules d'attention et MLP
(\texttt{q/k/v/o\_proj}, \texttt{gate/up/down\_proj}),
6~époques, $lr$\,=\,$8 \times 10^{-5}$, batch effectif 16
(batch\,=\,2, \texttt{gradient\_accumulation\_steps}\,=\,8),
optimiseur AdamW 8-bit, \texttt{max\_seq\_length}\,=\,1024.
Le format de prompting est Alpaca (\texttt{\#\#\# Instruction /
\#\#\# Input / \#\#\# Response}). L'inférence utilise un
\texttt{repetition\_penalty}\,=\,1.16 et
\texttt{no\_repeat\_ngram\_size}\,=\,4.

\subsection{Méthode 4 --- FT+RAG}

FT+RAG combine l'adaptateur issu du fine-tuning LoRA
(\texttt{lora\_adapter\_ft}) avec une récupération FAISS à l'inférence~:
les passages sont récupérés depuis le même index que le RAG (chunks
200/50) et injectés dans le prompt Alpaca avant décodage local.
Cette méthode est \textbf{distincte} du RAFT~: ici le retrieveur est
actif à l'inférence, non à l'entraînement. L'objectif est de tester
si la combinaison passages + adaptateur surpasse le RAG seul.

\subsection{Méthode 5 --- RAFT}

RAFT étend le fine-tuning LoRA en construisant les exemples
d'entraînement de manière à imiter un contexte RAG~: pour chaque
exemple de \texttt{train.json}, le retrieveur FAISS récupère les
top-$k$ passages, certains pertinents, d'autres servant de distracteurs.
Le modèle apprend ainsi à identifier et exploiter les passages utiles
tout en ignorant le bruit. L'adaptateur résultant
(\texttt{lora\_adapter\_raft}) est ensuite appliqué sur les questions
de test \textbf{sans} contexte injecté à l'inférence.

\subsection{Méthode 6 --- RAG + Reranking}

%%
%% [FIGURE 4 — À INSÉRER]
%% Pipeline RAG + Reranking (extension du schéma RAG)
%% Fichier : figures/fig4_rerank_pipeline.pdf
%%
\begin{figure}[H]
  \centering
  \fbox{\parbox{0.9\linewidth}{\centering
    \textit{[Figure 4 --- Pipeline RAG + Reranking]}\\[4pt]
    Question $\rightarrow$ FAISS (top-M) $\rightarrow$ Cross-encoder
    mMiniLMv2 $\rightarrow$ top-k rerankés $\rightarrow$ Groq
  }}
  \caption{Extension du RAG par reranking cross-encoder
           (\texttt{mmarco-mMiniLMv2-L12-H384-v1}).}
  \label{fig:rerank}
\end{figure}

Le même index FAISS est interrogé pour un top-M élargi de candidats,
puis un cross-encoder multilingue
(\texttt{cross-encoder/mmarco-mMiniLMv2-L12-H384-v1}) évalue chaque
paire (requête, passage) pour ne conserver que les top-$k$ passages
les plus pertinents, transmis à Groq pour génération.

\subsection{Méthode 7 --- Function Calling}

%%
%% [FIGURE 5 — À INSÉRER]
%% Diagramme de séquence du function calling (sequenceDiagram Mermaid)
%% Fichier : figures/fig5_function_calling.pdf
%%
\begin{figure}[H]
  \centering
  \fbox{\parbox{0.9\linewidth}{\centering
    \textit{[Figure 5 --- Diagramme de séquence Function Calling]}\\[4pt]
    Utilisateur $\rightarrow$ Groq (tour 1, schéma outil) $\rightarrow$
    \texttt{tool\_call search\_docs} $\rightarrow$ FAISS Python
    $\rightarrow$ Groq (tour 2, synthèse) $\rightarrow$ Réponse
  }}
  \caption{Pipeline function calling en deux tours. Le LLM décide
           quand déclencher la recherche documentaire.}
  \label{fig:fc}
\end{figure}

Le modèle reçoit la question et le schéma de l'outil \texttt{search\_docs}.
S'il décide de l'invoquer, l'exécution Python lance une recherche FAISS
et retourne les passages au modèle, qui génère ensuite la réponse finale.
Environ 12\,\% des requêtes ne déclenchent pas d'appel d'outil
(repli sans contexte).

\subsection{Métriques d'évaluation}
\label{subsec:metriques}

Nous mobilisons un panel de métriques complémentaires
(Table~\ref{tab:metriques}). Les intervalles de confiance à 95\,\% sont
estimés par \textbf{bootstrap} (1\,000 réplicatas, seed\,=\,42) sur
les scores par paire pour \bs{}, \met{}, F1 et \rl{}. Le proxy
d'hallucination est défini comme~:
\[
  \text{Hall} = 1 - \text{ROUGE-L}(\hat{y}, c)
\]
où $\hat{y}$ est la réponse prédite et $c$ le contexte source du jeu
de test. Le corpus étant 100\,\% français (HAL remplace Arxiv),
ce calcul est effectué dans la même langue pour tous les exemples,
éliminant le biais linguistique.

L'évaluation subjective est conduite via \textbf{LLM-as-judge}~:
Claude Sonnet~4.6 note chaque réponse selon trois critères
(cohérence, utilité, fidélité au contexte) sur une échelle 1--5.

\begin{table}[h]
  \centering
  \caption{Métriques d'évaluation utilisées}
  \label{tab:metriques}
  \renewcommand{\arraystretch}{1.2}
  \begin{tabular}{lll}
    \toprule
    \textbf{Métrique} & \textbf{Type} & \textbf{Référence} \\
    \midrule
    F1 token-level     & Lexicale    & \cite{rajpurkar2016squad} \\
    ROUGE-L            & Lexicale    & \cite{lin2004rouge} \\
    METEOR             & Lexicale    & \cite{banerjee2005meteor} \\
    BERTScore          & Sémantique  & \cite{zhang2020bertscore} \\
    IC 95\% bootstrap  & Statistique & --- \\
    Acc@BS (80/85/90\%) & Sémantique & --- \\
    Hallucination (proxy) & Fidélité & --- \\
    Faithfulness       & Fidélité    & --- \\
    LLM-as-judge       & Subjective  & --- \\
    Latence (ms)       & Efficacité  & --- \\
    CO$_2$eq (gCO$_2$eq) & Écologique & \cite{strubell2019energy} \\
    \bottomrule
  \end{tabular}
\end{table}

\subsection{Infrastructure et reproductibilité}

Les méthodes Baseline, RAG, Reranking et Function calling sont exécutées
via l'API Groq (modèle \texttt{llama-3.1-8b-instant}, H100 SXM partagé).
Le fine-tuning LoRA, FT+RAG et RAFT sont effectués sur GPU NVIDIA A100
(Google Colab), avec profil automatique vers T4/L4 si A100 indisponible.
Toutes les expériences utilisent seed\,=\,42. Le code est disponible à
l'adresse~: \url{https://github.com/coulby45/llm-integration-pipeline}.


%% ============================================================
\section{Résultats}
\label{sec:results}

\subsection{Comparaison globale}
\label{subsec:global}

Le tableau~\ref{tab:global} présente les résultats agrégés sur les
290 exemples de test. Le meilleur score par colonne est en
\textbf{\textcolor{best}{vert gras}} ; le moins bon en
\textcolor{worst}{rouge}.

\begin{table*}[t]
  \centering
  \caption{Résultats comparatifs des sept méthodes (n\,=\,290,
           campagne 2026-05-06). IC\,95\,\% bootstrap pour BERTScore,
           METEOR, F1 et ROUGE-L. $\downarrow$\,=\,plus bas est meilleur.}
  \label{tab:global}
  \renewcommand{\arraystretch}{1.3}
  \resizebox{\textwidth}{!}{%
  \begin{tabular}{lcccccccccc}
    \toprule
    \multirow{2}{*}{\textbf{Méthode}}
      & \textbf{F1}     & \textbf{\bs}   & \textbf{IC 95\% BS}
      & \textbf{\rl}    & \textbf{\met}
      & \textbf{Fidélit.} & \textbf{Hall.}$\downarrow$
      & \textbf{Acc@85\%} & \textbf{Latence}$\downarrow$ \\
      & (\%)            & (\%)           & (\%)
      & (\%)            & (\%)
      & (\%)              & (\%)
      & (\%)              & (ms) \\
    \midrule
    Baseline
      & 16.3            & 81.4           & [81.11,\,81.72]
      & 13.7            & 34.5
      & 16.5            & 83.5
      & 9.7             & \best{549} \\
    RAG
      & 32.8            & 85.8           & [85.22,\,86.23]
      & 28.7            & 52.4
      & 31.8            & 68.2
      & 54.5            & 599 \\
    Fine-tuné (LoRA)
      & \worst{14.8}    & 81.4           & [81.08,\,81.64]
      & 11.6            & 35.5
      & 12.7            & 87.3
      & 6.2             & 14\,217 \\
    FT+RAG
      & \worst{14.6}    & \worst{81.0}   & [80.69,\,81.27]
      & 11.7            & 35.8
      & 13.1            & 86.9
      & 5.9             & 18\,976 \\
    RAFT
      & \worst{11.9}    & 81.2           & [80.94,\,81.54]
      & \worst{10.1}    & \worst{26.3}
      & \worst{11.6}    & \worst{88.4}
      & \worst{4.8}     & 21\,224 \\
    \textbf{Rerank}
      & \best{36.5}     & \best{86.6}    & \best{[86.05,\,87.12]}
      & \best{32.1}     & \best{56.1}
      & \best{33.8}     & \best{66.2}
      & \best{60.7}     & 570 \\
    Function calling
      & 23.9            & 83.4           & [82.86,\,83.89]
      & 20.9            & 44.6
      & 24.2            & 75.8
      & 31.4            & 1\,078 \\
    \bottomrule
  \end{tabular}}
\end{table*}

%%
%% [FIGURE 6 — À INSÉRER]
%% Barres groupées des métriques globales
%% Fichier : results/plots/fig_metriques_globales.png
%%
\begin{figure}[H]
  \centering
  \fbox{\parbox{0.9\linewidth}{\centering
    \textit{[Figure 6 --- Métriques globales --- barres groupées]}\\[4pt]
    F1, BERTScore, ROUGE-L, METEOR, Fidélité par méthode\\
    Source~: \texttt{results/plots/fig\_metriques\_globales.png}
  }}
  \caption{Comparaison des sept méthodes sur les métriques globales.}
  \label{fig:barres_globales}
\end{figure}

Le \textbf{reranking} arrive en tête sur toutes les métriques clés~:
F1 (36,5\,\%), \bs{} (86,6\,\%), \rl{} (32,1\,\%) et \met{}
(56,1\,\%), suivi du \textbf{RAG} (F1\,=\,32,8\,\%). Le
\textit{function calling} se positionne en intermédiaire (F1\,=\,23,9\,\%).
Les méthodes de fine-tuning (LoRA, FT+RAG, RAFT) \textbf{n'améliorent
pas} la baseline sur la majorité des métriques~: le RAFT obtient le
F1 le plus bas (11,9\,\%), et les trois méthodes GPU affichent un
proxy d'hallucination supérieur à la baseline (83,5\,\%).

Concernant la latence, les méthodes GPU (LoRA~: 14\,217\,ms,
FT+RAG~: 18\,976\,ms, RAFT~: 21\,224\,ms) sont respectivement
\textbf{26$\times$}, \textbf{35$\times$} et \textbf{39$\times$}
plus lentes que la baseline (549\,ms), tandis que le reranking (570\,ms)
n'induit qu'un surcoût marginal de 4\,\%.

\subsection{Analyse par domaine}
\label{subsec:domaine}

%%
%% [FIGURE 7 — À INSÉRER]
%% Heatmap BERTScore méthodes × dataset_type
%% Fichier : results/plots/fig_heatmap_bertscore.png
%%
\begin{figure}[H]
  \centering
  \fbox{\parbox{0.9\linewidth}{\centering
    \textit{[Figure 7 --- Heatmap BERTScore méthodes $\times$ domaines]}\\[4pt]
    Source~: \texttt{results/plots/fig\_heatmap\_bertscore.png}
  }}
  \caption{BERTScore (\%) par méthode et par domaine.}
  \label{fig:heatmap}
\end{figure}

\begin{table}[h]
  \centering
  \caption{BERTScore (\%) par domaine et par méthode}
  \label{tab:domaines}
  \renewcommand{\arraystretch}{1.2}
  \begin{tabular}{lccccc}
    \toprule
    \textbf{Méthode} & \textbf{Tech.} & \textbf{Multi-s.}
      & \textbf{Temp.} & \textbf{Jurid.} \\
    \midrule
    Baseline        & 82.0 & 81.8 & 80.5 & 81.7 \\
    RAG             & 85.9 & 86.9 & 84.9 & 84.5 \\
    Fine-tuné       & 82.0 & 81.4 & 81.1 & 80.2 \\
    FT+RAG          & 80.8 & 81.1 & 81.2 & 80.2 \\
    RAFT            & 81.3 & 81.6 & 81.1 & 80.0 \\
    \textbf{Rerank} & \best{86.6} & \best{87.2}
                    & \best{86.3} & \best{85.2} \\
    Funct. call.    & 83.5 & 83.7 & 83.0 & 83.6 \\
    \bottomrule
  \end{tabular}
\end{table}

Le reranking excelle sur tous les domaines, avec un pic en
multi-sauts (87,2\,\%). Le domaine temporel révèle la plus forte
amélioration RAG sur baseline (+4,4\,pts de BERTScore), confirmant
l'intérêt de l'injection documentaire pour les contenus post-coupure.
Le domaine juridique (code de la route) reste difficile pour toutes
les méthodes hors retrieval~: le RAFT y atteint son score d'hallucination
le plus élevé (92,8\,\%, Table~\ref{tab:hallucination}).

\subsection{Analyse par type de question}

\begin{table}[h]
  \centering
  \caption{BERTScore (\%) par type de question}
  \label{tab:qtypes}
  \renewcommand{\arraystretch}{1.2}
  \begin{tabular}{lccccc}
    \toprule
    \textbf{Méthode} & \textbf{Fact.} & \textbf{Synth.}
      & \textbf{Compr.} & \textbf{Simple} & \textbf{Compl.} \\
    \midrule
    Baseline    & 82.0 & 80.7 & 80.9 & 81.9 & 81.7 \\
    RAG         & 86.4 & 83.7 & 85.8 & 88.5 & 84.9 \\
    Fine-tuné   & 81.3 & 81.6 & 80.9 & 81.0 & 81.8 \\
    FT+RAG      & 80.5 & 81.3 & 81.2 & 80.2 & 82.1 \\
    RAFT        & 81.0 & 81.4 & 80.9 & 82.1 & 80.9 \\
    Rerank      & \best{88.5} & \best{84.4} & \best{85.5}
                & \best{89.0} & \best{85.0} \\
    Funct. call.& 84.5 & 82.0 & 83.2 & 83.9 & 83.4 \\
    \bottomrule
  \end{tabular}
\end{table}

%%
%% [FIGURE 8 — À INSÉRER]
%% Comparaison questions simples vs complexes HAL
%% Fichier : results/plots/fig_multisauts.png
%%
\begin{figure}[H]
  \centering
  \fbox{\parbox{0.9\linewidth}{\centering
    \textit{[Figure 8 --- Questions simples vs multi-sauts (HAL)]}\\[4pt]
    BERTScore simple vs complexe par méthode\\
    Source~: \texttt{results/plots/fig\_multisauts.png}
  }}
  \caption{Impact de la complexité des questions sur le BERTScore
           (corpus HAL).}
  \label{fig:multisauts}
\end{figure}

Un résultat notable concerne les questions \textbf{simples} (HAL)~:
le RAG (88,5\,\%) et le reranking (89,0\,\%) y obtiennent leurs
meilleurs scores absolus. Le RAG surpasse le reranking sur les
questions multi-sauts complexes de 0,1\,pt (84,9\,\% vs 85,0\,\%),
une différence négligeable. Le \textbf{FT+RAG} inverse ce résultat~:
il performe légèrement mieux sur les complexes (82,1\,\%) que sur
les simples (80,2\,\%), tandis que le \textbf{RAFT} montre la
tendance inverse (82,1\,\% simple vs 80,9\,\% complexe), suggérant
que l'entraînement avec distracteurs n'a pas conduit au comportement
multi-sauts escompté.

\subsection{Taux d'hallucination}

%%
%% [FIGURE 9 — À INSÉRER]
%% Taux d'hallucination global + par domaine
%% Fichier : results/plots/fig_hallucination.png
%%
\begin{figure}[H]
  \centering
  \fbox{\parbox{0.9\linewidth}{\centering
    \textit{[Figure 9 --- Taux d'hallucination par domaine]}\\[4pt]
    Source~: \texttt{results/plots/fig\_hallucination.png}
  }}
  \caption{Proxy d'hallucination ($1 - \text{ROUGE-L}(\hat{y}, c)$) par
           méthode et par domaine (\%). Plus bas = moins d'hallucination.}
  \label{fig:hallucination}
\end{figure}

\begin{table}[h]
  \centering
  \caption{Taux d'hallucination proxy (\%) par domaine. $\downarrow$ = meilleur.}
  \label{tab:hallucination}
  \renewcommand{\arraystretch}{1.2}
  \begin{tabular}{lcccc}
    \toprule
    \textbf{Méthode} & \textbf{Tech.} & \textbf{Temp.}
      & \textbf{Multi-s.} & \textbf{Jurid.} \\
    \midrule
    Baseline    & 80.2 & 86.4 & 84.2 & 79.5 \\
    RAG         & 67.5 & 73.4 & 63.7 & 66.1 \\
    Fine-tuné   & 85.7 & 87.7 & 88.0 & 88.6 \\
    FT+RAG      & 86.5 & 86.8 & 87.2 & 87.9 \\
    RAFT        & 88.4 & 88.6 & 87.2 & \worst{92.8} \\
    \textbf{Rerank}
                & \best{63.0} & \best{71.5}
                & \best{63.5} & \best{64.3} \\
    Funct. call.& 73.9 & 78.5 & 76.2 & 69.0 \\
    \bottomrule
  \end{tabular}
\end{table}

Le reranking atteint le proxy d'hallucination le plus bas sur tous les
domaines, avec 63,0\,\% en technique et 64,3\,\% en juridique. Le
fine-tuning LoRA (87,3\,\%), FT+RAG (86,9\,\%) et RAFT (88,4\,\%)
présentent des taux d'hallucination \textbf{supérieurs à la baseline}
(83,5\,\%), ce qui indique que les réponses générées localement
s'écartent davantage du contexte de référence.

\subsection{Précision par seuil BERTScore}

%%
%% [FIGURE 10 — À INSÉRER]
%% Accuracy@BS par seuil (80/85/90%)
%% Fichier : results/plots/fig_accuracy_bs.png
%%
\begin{figure}[H]
  \centering
  \fbox{\parbox{0.9\linewidth}{\centering
    \textit{[Figure 10 --- Accuracy@BERTScore par seuil]}\\[4pt]
    Acc@80\%, Acc@85\%, Acc@90\% par méthode\\
    Source~: \texttt{results/plots/fig\_accuracy\_bs.png}
  }}
  \caption{Pourcentage de réponses atteignant un seuil de BERTScore
           de 80\,\%, 85\,\% et 90\,\%.}
  \label{fig:acc_bs}
\end{figure}

À seuil 85\,\%, le reranking atteint 60,7\,\% de réponses qualifiées,
le RAG 54,5\,\%, contre seulement 9,7\,\% pour la baseline et 5,9\,\%
pour le FT+RAG. Aucune méthode ne dépasse 21\,\% au seuil strict de
90\,\%, illustrant la difficulté intrinsèque du jeu de test français.

\subsection{Évaluation subjective LLM-as-judge}
\label{subsec:judge}

Une évaluation complémentaire a été conduite via Claude Sonnet~4.6 en
tant que juge automatique, notant chaque réponse sur trois critères
(cohérence, utilité, fidélité) de 1 à 5.

\begin{table}[h]
  \centering
  \caption{Scores LLM-as-judge (Claude Sonnet 4.6, n\,=\,290).
           Composite = moyenne des trois critères.}
  \label{tab:judge}
  \renewcommand{\arraystretch}{1.2}
  \begin{tabular}{lcccc}
    \toprule
    \textbf{Méthode} & \textbf{Cohér.} & \textbf{Util.}
      & \textbf{Fidél.} & \textbf{Compos.} \\
    \midrule
    Baseline       & 2.46 & 1.72 & 1.37 & 1.85 \\
    RAG            & 3.68 & 3.31 & 3.29 & 3.43 \\
    Fine-tuné      & 2.02 & 1.61 & 1.25 & 1.63 \\
    FT+RAG         & 2.00 & 1.79 & 1.55 & 1.78 \\
    RAFT           & 1.63 & 1.53 & 1.28 & 1.48 \\
    \textbf{Rerank}& \best{3.93} & \best{3.66} & \best{3.57} & \best{3.72} \\
    Funct. call.   & 3.21 & 2.74 & 2.48 & 2.81 \\
    \bottomrule
  \end{tabular}
\end{table}

Le classement subjectif reproduit fidèlement l'ordre des métriques
automatiques~: rerank $>$ RAG $>$ function calling $>$ baseline $\approx$
FT+RAG $>$ fine-tuné $>$ RAFT. Le FT+RAG obtient un score composite
(1,78) comparable à la baseline (1,85), confirmant que la combinaison
LoRA + retrieval n'apporte pas de valeur ajoutée subjective sur ce
protocole.

\subsection{Impact écologique}

%%
%% [FIGURE 11 — À INSÉRER]
%% Impact écologique — énergie, CO2
%% Fichier : results/plots/fig_ecologie.png
%%
\begin{figure}[H]
  \centering
  \fbox{\parbox{0.9\linewidth}{\centering
    \textit{[Figure 11 --- Impact écologique]}\\[4pt]
    Énergie (Wh) et CO$_2$ France / mondial par méthode\\
    Source~: \texttt{results/plots/fig\_ecologie.png}
  }}
  \caption{Estimation de l'empreinte énergétique et carbone.
           Formule~: $E = \text{latence}_s \times \text{TDP} \times
           \text{PUE} / 3600$, avec intensité carbone France
           52\,gCO$_2$/kWh.}
  \label{fig:ecologie}
\end{figure}

\begin{table}[h]
  \centering
  \caption{Empreinte énergétique estimée pour le jeu de test complet
           (n\,=\,290). TDP API~: 50\,W, PUE 1,10~;
           TDP GPU local (A100)~: 150\,W, PUE 1,15.}
  \label{tab:ecologie}
  \renewcommand{\arraystretch}{1.2}
  \begin{tabular}{lccc}
    \toprule
    \textbf{Méthode} & \textbf{Énergie (Wh)} &
    \textbf{CO$_2$ FR (g)} & \textbf{Ratio vs Baseline} \\
    \midrule
    Baseline         & 2.43  & 0.13  & $\times$1.0  \\
    \textbf{Rerank}  & 2.52  & 0.13  & $\times$1.04 \\
    RAG              & 2.65  & 0.14  & $\times$1.09 \\
    Function calling & 4.77  & 0.25  & $\times$2.0  \\
    Fine-tuné        & 197.6 & 10.3  & $\times$81   \\
    FT+RAG           & 263.7 & 13.7  & $\times$109  \\
    RAFT             & 294.9 & 15.3  & $\times$121  \\
    \bottomrule
  \end{tabular}
\end{table}

Les méthodes GPU (LoRA, FT+RAG, RAFT) concentrent l'essentiel de
l'empreinte énergétique~: le RAFT consomme \textbf{121 fois} plus
que la baseline par prédiction. Le reranking, grâce à sa latence
minimale (570\,ms), présente la meilleure efficacité énergétique
parmi les méthodes avec contexte, avec un surcoût de seulement
4\,\% par rapport à la baseline.


%% ============================================================
\section{Discussion}
\label{sec:discussion}

\subsection{RAG et reranking comme stratégie dominante}

Les résultats convergent clairement vers la supériorité des méthodes
d'injection documentaire à l'inférence. Le reranking améliore le RAG
seul de 3,7\,pts de F1 pour un surcoût de latence négligeable
($+$4\,\%). Ce gain s'explique par la qualité du passage sélectionné~:
le cross-encoder reclasse les candidats selon la pertinence de la
paire (question, passage), corrigeant les erreurs du retrieveur dense.
La qualité du corpus indexé reste le facteur déterminant~: un
paramétrage de chunks trop courts (50--100 mots) dégradait le RAG
sous la baseline dans une configuration préliminaire~; les chunks
de 200 mots avec chevauchement de 50 mots représentent le meilleur
compromis identifié.

\subsection{Échec relatif du fine-tuning sur ce protocole}

Le fait que LoRA, FT+RAG et RAFT n'améliorent pas la baseline
s'explique par plusieurs facteurs concomitants~:

\paragraph{Désalignement de quantification.}
L'inférence Groq (\texttt{llama-3.1-8b-instant}) et l'inférence
locale (4-bit Unsloth, A100) produisent des distributions de sortie
distinctes, même à paramètres identiques.

\paragraph{Désalignement de gabarit.}
Le format Alpaca utilisé à l'entraînement (\texttt{\#\#\# Instruction
/ Response}) diffère du prompt libre employé à l'inférence de test,
créant une distorsion entre les distributions d'entraînement et
d'évaluation. Pour FT+RAG, l'injection de passages au moment de
l'inférence ne correspond pas exactement au gabarit employé en SFT,
ce qui peut expliquer l'absence de synergie observée.

Ces deux sources de biais doivent être prises en compte avant de conclure à
l'infériorité intrinsèque du fine-tuning.

\subsection{Function calling~: intérêt agentique, coût de fiabilité}

Le function calling offre une flexibilité architecturale importante
(recherche conditionnelle, intégration dans des workflows agentiques)
mais présente une qualité inférieure au RAG direct dans ce benchmark.
Le taux de repli sans outil ($\approx$\,12\,\%) signifie qu'une fraction
non négligeable des réponses est générée sans contexte documentaire.
Sa latence ($\times$2 vs baseline) reste acceptable pour un déploiement
agentique.

\subsection{Limites méthodologiques}
\label{subsec:limites}

\paragraph{Taille du corpus.}
Avec 290 exemples de test répartis sur quatre domaines, certains
croisements méthode $\times$ domaine reposent sur moins de 30 paires
(notamment le juridique, $n_\text{test}$\,=\,20), limitant la puissance
statistique des comparaisons locales.

\paragraph{Évaluation automatique.}
Les métriques lexicales et \bs{} pénalisent les paraphrases correctes
et ne capturent pas entièrement la cohérence factuelle. L'évaluation
LLM-as-judge constitue une première réponse à cette limite, mais reste
elle-même dépendante du modèle juge utilisé.

\paragraph{Proxy d'hallucination.}
Le proxy $1 - \text{ROUGE-L}(\hat{y}, c)$ mesure l'écart lexical
entre réponse prédite et contexte source, et non la vérité factuelle.
Une réponse correcte mais formulée différemment sera pénalisée.

\subsection{Recommandations pratiques}

En synthèse, nos recommandations pour un déploiement en contexte
francophone contraint sont les suivantes. Lorsque la priorité est la
qualité et la fraîcheur documentaire, RAG + reranking est la meilleure
option~: le surcoût de latence est marginal ($+$4\,\% vs baseline) et
le gain en F1 est de $+$20\,pts. Sous contrainte GPU nulle, le RAG
seul offre un excellent compromis avec une mise à jour instantanée
par réindexation. Pour un flux agentique ou une politique de recherche
conditionnelle, le function calling est adapté, au prix d'une qualité
légèrement inférieure. Enfin, le fine-tuning LoRA n'est à envisager
que si le domaine cible est stable, très spécialisé, et si le gabarit
de prompt est rigoureusement aligné entre entraînement et inférence.


%% ============================================================
\section{Conclusion}
\label{sec:conclusion}

Nous avons présenté une comparaison systématique de sept stratégies
d'intégration de connaissances dans \textsc{LLaMA}~3.1~8B, sur un
corpus 100\,\% français de 290 paires question-réponse couvrant quatre
domaines (technique, multi-sauts HAL, temporel, juridique).

Nos résultats montrent que le \textbf{reranking} et le \textbf{RAG}
constituent les approches les plus efficaces, réduisant le taux
d'hallucination de plus de 17\,pts et améliorant le F1 d'un facteur
supérieur à 2 par rapport à la baseline. Ces conclusions sont
cohérentes entre métriques automatiques et évaluation subjective
LLM-as-judge. Les méthodes de fine-tuning (LoRA, FT+RAG, RAFT)
n'ont pas amélioré la baseline dans ce protocole, principalement
en raison du désalignement de quantification entre inférence API
et locale, ainsi que des différences de gabarit de prompt entre
entraînement et évaluation.
Le function calling se positionne comme une alternative intéressante
pour les architectures agentiques. Sur le plan écologique, les
méthodes GPU concentrent l'essentiel de l'empreinte carbone, avec
un ratio allant jusqu'à $\times$121 pour le RAFT vis-à-vis de la
baseline API.

\paragraph{Perspectives.}
Plusieurs pistes méritent d'être explorées~: (i)~harmonisation du
gabarit de prompt entre SFT et inférence pour FT+RAG~;
(ii)~ablations retrieveur (BM25 vs dense, fusion hybride)~;
(iii)~mesure directe du taux d'erreurs JSON pour le function
calling~; (iv)~évaluation humaine pour calibrer les métriques
automatiques~; (v)~retrieveurs plus riches (DPR, ColBERT) pour
les questions multi-sauts.

\paragraph{Code et reproductibilité.}
L'ensemble du code (notebooks \texttt{01}--\texttt{09}) est disponible
sur GitHub~: \url{https://github.com/coulby45/llm-integration-pipeline}.
Les graines aléatoires (seed\,=\,42), configurations matérielles
(A100 / Groq H100 partagé) et versions de librairies sont documentées
dans \texttt{requirements.txt}.


%% ============================================================
\section*{Remerciements}

Ce travail a été réalisé dans le cadre d'un projet pédagogique à
CESI École d'Ingénieurs, Saint-Nazaire. L'auteur remercie les équipes
de Groq pour l'accès API, ainsi que les contributeurs d'Unsloth pour
la librairie de fine-tuning efficace.


%% ============================================================
\begin{thebibliography}{99}

\bibitem{vaswani2017attention}
A. Vaswani \textit{et al.},
``Attention is all you need,''
in \textit{Advances in Neural Information Processing Systems (NeurIPS)},
2017.

\bibitem{brown2020language}
T. B. Brown \textit{et al.},
``Language models are few-shot learners,''
in \textit{NeurIPS}, 2020.

\bibitem{radford2019language}
A. Radford \textit{et al.},
``Language models are unsupervised multitask learners,''
\textit{OpenAI Blog}, 2019.

\bibitem{touvron2023llama}
H. Touvron \textit{et al.},
``LLaMA 2: Open foundation and fine-tuned chat models,''
\textit{arXiv:2307.09288}, 2023.

\bibitem{meta2024llama3}
Meta AI,
``Meta LLaMA 3,''
\textit{arXiv:2407.21783}, 2024.

\bibitem{jiang2023mistral}
A. Q. Jiang \textit{et al.},
``Mistral 7B,''
\textit{arXiv:2310.06825}, 2023.

\bibitem{openai2023gpt4}
OpenAI,
``GPT-4 Technical Report,''
\textit{arXiv:2303.08774}, 2023.

\bibitem{lewis2020rag}
P. Lewis \textit{et al.},
``Retrieval-augmented generation for knowledge-intensive NLP tasks,''
in \textit{NeurIPS}, 2020.

\bibitem{johnson2019faiss}
J. Johnson, M. Douze, and H. Jégou,
``Billion-scale similarity search with GPUs,''
\textit{IEEE Trans. Big Data}, vol. 7, no. 3, 2021.

\bibitem{reimers2019sentencebert}
N. Reimers and I. Gurevych,
``Sentence-BERT: Sentence embeddings using Siamese BERT-networks,''
in \textit{EMNLP}, 2019.

\bibitem{nogueira2019passage}
R. Nogueira and K. Cho,
``Passage Re-ranking with BERT,''
\textit{arXiv:1901.04085}, 2019.

\bibitem{hu2022lora}
E. J. Hu \textit{et al.},
``LoRA: Low-rank adaptation of large language models,''
in \textit{ICLR}, 2022.

\bibitem{dettmers2023qlora}
T. Dettmers \textit{et al.},
``QLoRA: Efficient finetuning of quantized LLMs,''
in \textit{NeurIPS}, 2023.

\bibitem{hong2024raft}
T. Zhang \textit{et al.},
``RAFT: Adapting language model to domain specific RAG,''
\textit{arXiv:2403.10131}, 2024.

\bibitem{schick2023toolformer}
T. Schick \textit{et al.},
``Toolformer: Language models can teach themselves to use tools,''
in \textit{NeurIPS}, 2023.

\bibitem{ji2023hallucination}
Z. Ji \textit{et al.},
``Survey of hallucination in natural language generation,''
\textit{ACM Computing Surveys}, vol. 55, no. 12, 2023.

\bibitem{rajpurkar2016squad}
P. Rajpurkar \textit{et al.},
``SQuAD: 100,000+ questions for machine comprehension of text,''
in \textit{EMNLP}, 2016.

\bibitem{lin2004rouge}
C.-Y. Lin,
``ROUGE: A package for automatic evaluation of summaries,''
in \textit{ACL Workshop on Text Summarization}, 2004.

\bibitem{banerjee2005meteor}
S. Banerjee and A. Lavie,
``METEOR: An automatic metric for MT evaluation with improved
correlation with human judgments,''
in \textit{ACL Workshop}, 2005.

\bibitem{zhang2020bertscore}
T. Zhang \textit{et al.},
``BERTScore: Evaluating text generation with BERT,''
in \textit{ICLR}, 2020.

\bibitem{es2024ragas}
S. Es \textit{et al.},
``RAGAS: Automated evaluation of retrieval augmented generation,''
in \textit{EACL Demo Track}, 2024.

\bibitem{mccloskey1989catastrophic}
M. McCloskey and N. J. Cohen,
``Catastrophic interference in connectionist networks,''
\textit{Psychology of Learning and Motivation}, vol. 24, 1989.

\bibitem{kirkpatrick2017overcoming}
J. Kirkpatrick \textit{et al.},
``Overcoming catastrophic forgetting in neural networks,''
\textit{Proceedings of the National Academy of Sciences}, 2017.

\bibitem{meng2022locating}
K. Meng \textit{et al.},
``Locating and editing factual associations in GPT,''
in \textit{NeurIPS}, 2022.

\bibitem{mitchell2022memory}
E. Mitchell \textit{et al.},
``Memory-based model editing at scale,''
in \textit{ICML}, 2022.

\bibitem{strubell2019energy}
E. Strubell, A. Ganesh, and A. McCallum,
``Energy and policy considerations for deep learning in NLP,''
in \textit{ACL}, 2019.

\bibitem{unsloth2024}
D. Han and M. Han,
``Unsloth: Fast and memory-efficient LLM fine-tuning,''
\textit{GitHub Repository}, 2024.
\url{https://github.com/unslothai/unsloth}

\end{thebibliography}

\end{document}